% Appendix B: the answers to every part of every exercise that is not code, one section per
% chapter. The code parts are left out on purpose: each exercise names the testbench or host suite
% that checks it, and that check is the authority. Where an answer says what a testbench or a suite
% reports, it quotes what a design built to the chapters' specifications produced when it was run.

\chapter{Answers to the Exercises}
\label{sol:ch}

{\large\itshape Every part that is not code, answered in full.\par}

These are the answers to the parts of the exercises that you answer in words or numbers: the
reasoning, the calculations, the traces, and the parts that ask you to break a working design on
purpose and say what happens. They are in the order the exercises are set, one section per chapter
that has any. Chapters~6 to~8 have none: every exercise there is a header or a class, and the host
suite of \chapref{8} is its check.
Every answer opens with the exercise it answers, linked back to it and to the page it is on, and each part's
answer is lettered as the part is. A part that is not listed is code, and has nothing to answer
here.

\textbf{The VHDL and C++ the exercises ask for are not answered.} Each exercise names the testbench
or the host suite that checks it, and that check is the authority on your module: it is the
executable form of the specification the chapter sets, and a module that passes it has met that
specification, however it is written. A listing printed here would be one answer of many, and it
would teach less than the testbench's failure message does.

Where a part asks what a testbench or a suite reports when you break something, the answer quotes
the message a design built to the chapter's specification actually produced when it was broken in
exactly that way. If yours reports something else, the difference is worth understanding before you
move on: it usually means your design differs from the specification somewhere the testbench does
not look.

\textbf{Try each part before you read its answer.} That matters most for the parts that ask you to
predict: an outcome read before it is predicted is an outcome you can no longer disagree with, and
the disagreement is the lesson.

% Section B.1: the answers to the parts of lectures/L01/appendix/c_exercises.md that are not code.
% Every testbench message quoted was produced by a design built to Chapters 1 to 5 and broken in
% exactly the way the part describes.

\section{Chapter 1: The Register Package and the Peripheral Top}
\label{sol:c1}
Exercise~\ref{c1:ex:2} is the build of \code{uart_top}, checked by a clean analysis now and by
\code{uart_top_tb} from \chapref{5} on; only its closing question, part e), is answered here.

\solution{c1:ex:1}{\code{uart_def}}
\ans{a} The register indices are \code{REG_STATUS} 0, \code{REG_CTRL} 1, \code{REG_BAUD_DIV} 2,
\code{REG_TX_DATA} 3, \code{REG_RX_DATA} 4, \code{REG_RX_POP} 5 and \code{REG_ERR_FLAGS} 6, each the
byte offset divided by four. The \code{STATUS} bits are \code{ST_TX_READY} 0, \code{ST_RX_VALID} 1,
\code{ST_ERROR} 2 and \code{ST_TX_IDLE} 3; the \code{CTRL} bits \code{CT_ENABLE} 0,
\code{CT_PARITY_LO} 1, \code{CT_PARITY_HI} 2, \code{CT_STOP} 3, \code{CT_RX_IRQ} 4 and
\code{CT_TX_IRQ} 5; the \code{ERROR_FLAGS} bits \code{ER_FRAMING} 0, \code{ER_PARITY} 1 and
\code{ER_OVERRUN} 2.

\ans{b} In VHDL the position indexes the vector directly:

\begin{vhdlcode}
if (status(ST_RX_VALID) = '1') then
\end{vhdlcode}

\noindent In C++ the same number is shifted into a mask at the use site:
\code{(status & (1U << status::RX_VALID)) != 0U}. The protocol speaks of ``\code{STATUS} bit 1'',
and a position is that number and nothing else, so both files can be checked against the protocol
by reading it off, digit for digit. A mask would be a derived value, \code{2} here and \code{4} for
bit 2, which the reader has to recompute to compare, and which only means something for one word
width. Storing positions also keeps the two languages identical, because the only difference between
them is the operator that uses the constant, and that operator is written where it is used.

\ans{c} \textbf{In no testbench at all.} Every testbench that touches those two registers names them
through \code{uart_def}, and so does the register bank. \code{uart_regs_tb} writes \code{RX_POP} to
``\code{REG_RX_POP}'' and the bank decodes ``\code{REG_RX_POP}'', so after the swap both sides mean
6 and the bench still passes; the same holds for \code{ERROR_FLAGS} at 5. \code{uart_top_tb} never
reads or writes index 5 or 6. Both were run with the swapped package, and both report that they
passed. A constant that is wrong in the one place everything reads it from is wrong consistently,
and a consistent error is invisible to any check built from the same source.

The C++ side is a separate transcription, and it still says \code{RX_POP} is 5 and
\code{ERROR_FLAGS} is 6, so the host suites in Chapters~8 and~9 pass too. The swap surfaces only
when the two transcriptions finally meet, on the bench in \chapref{10}. There the driver's
\code{read()} ends with \code{writeReg(RX_POP, 1)}, which the peripheral now takes as an
\code{ERROR_FLAGS} write: it sets the framing flag, so \code{STATUS} bit 2 comes up, and it never
advances the RX FIFO, so \code{RX_VALID} stays set and every later \code{read()} returns the same
byte. \code{clearErrors()} writes index 6, which the peripheral now takes as a pop, and silently
throws a received byte away.

Rung c of the ladder can still pass, because it sends one byte and reads one byte back, and that
first read is correct; it shows the fault only if the bring-up program looks at \code{STATUS} after
the read, or sends a second byte. Rung d shows it for certain: the first character typed comes back
correctly and so does every read after it, the same character each time. At rung e,
\code{EchoNode} echoes the first character typed without end.

\solution{c1:ex:2}{\code{uart_top}}
\ans{e} With only the provided files and your skeleton in \file{hw/}, \code{make build-vhdl} analyzes
five modules, runs the two transport testbenches, and skips the rest:

\begin{console}
==> uart_top_tb (skipped: no baud_gen.vhd sync.vhd fifo.vhd uart_tx.vhd uart_rx.vhd uart_regs.vhd yet)
2 testbench(es) run, 7 skipped.
\end{console}

\noindent \code{baud_gen} and \code{uart_tx} arrive in \chapref{2}, \code{sync} and \code{uart_rx} in
\chapref{3}, \code{fifo} in \chapref{4} and \code{uart_regs} in \chapref{5}. The one that
\code{uart_top} never instantiates is \code{fifo}: \code{uart_regs} owns both FIFOs. The build needs
it anyway because elaboration needs every entity in the hierarchy, not only the top's direct
children, and \code{uart_regs} cannot be analyzed until the entity it instantiates twice is already
in the \code{work} library.

\solution{c1:ex:3}{What a clean analysis does not prove}
\ans{a} \code{ghdl -a} accepts it. Both ports are \code{in std_logic}, so the entity is as legal
with them swapped as without, and the architecture refers to them by name, so it still analyzes.
Nothing binds to the entity positionally until \code{uart_top_tb} runs in \chapref{5}, and there
the first check fails: the bench's clock now arrives on the peripheral's \code{mosi} and its data on
\code{sclk}, so the first write never lands and the read-back reports

\begin{console}
uart_top_tb: BAUD_DIV read-back mismatch, got 0x0000!
\end{console}

\noindent four chapters after the mistake was made, in a message about \code{BAUD_DIV}. That is the
cost of a positional contract: the order of the ports is checked by nothing except a simulation of
the whole system, and when that simulation fails it points at the first symptom rather than the
cause. A named association would have caught it at analysis; the course binds positionally so that
the order in each interface table is something you have to get right.

\ans{b} All zeros. Every read transaction clocks out \code{0x00000000}, whatever register it
addresses, because the bridge latches \code{reg_rdata} one cycle after the command byte and shifts
it out on \code{MISO}. It is harmless until \chapref{5} because nothing reads it: the only testbench
that drives \code{uart_top} is skipped until then, and no register exists yet to be misreported.
Without the placeholder the vector is undriven, so every bit is \code{'U'}, and the master samples
\code{'U'} on all 32 data bits of every read. Both were simulated: the skeleton with the placeholder
returned \code{00000000000000000000000000000000} for \code{BAUD_DIV} and for \code{STATUS}, and the
skeleton without it returned 32 \code{'U'}s for each. The placeholder makes the answer a defined
value; it does not make it a right one.

\ans{c} A release that lands inside the recovery and removal window of the flip-flops around a
clock edge leaves some of them out of reset on that edge and others on the next, and any of them
may go metastable on the way. The design then starts half-reset: a state machine whose state
register has left reset while its outputs have not, a FIFO whose pointers start one cycle apart,
the SPI slave's edge detector seeing a false edge on a synchronizer stage that left reset a cycle
early. None of those states is one the logic was designed to leave, and nothing brings it back
short of another reset.

\code{sync} cannot do this job because of its port list. It takes \code{reset_s2_n} as an input:
it needs a clean reset in order to reset itself, so it cannot be the thing that produces one. A reset
synchronizer needs the raw \code{reset_n} on its flops' asynchronous clear, so that the assertion
reaches them with no clock at all, and a constant \code{'1'} clocked through the chain, so that the
release arrives a clean two edges later. \code{sync} has no port for either: its only asynchronous
input is the already-synchronized reset, and its data input is what it synchronizes.

\ans{d} Because the top's instantiation line is the block's contract. \code{uart_top} declares the
signals it connects and binds them by position, so writing that line fixes how many ports the
block has, in which order, of which types and in which directions, before the block exists. For
\code{baud_gen} to bind without the top changing, its entity must declare exactly
\code{clock} and \code{reset_s2_n} (\code{in std_logic}), then \code{div}
(\code{in natural range 1 to 65535}), then \code{tick} (\code{out std_logic}), in that order, which is
also the order \code{baud_gen_tb} binds. The names inside may be anything; the positions and the
types may not.

% Section B.2: the answers to the parts of lectures/L02/appendix/c_exercises.md that are not code.
% Every testbench message quoted was produced by a design built to Chapters 1 to 5 and broken in
% exactly the way the part describes.

\section{Chapter 2: UART Framing and the Transmitter}
\label{sol:c2}
The modules themselves are code: \code{baud_gen} and \code{uart_tx} in the a) parts of
Exercises~\ref{c2:ex:1} and~\ref{c2:ex:2}, checked by \code{baud_gen_tb} and \code{uart_tx_tb};
their instantiation into \code{uart_top} in Exercise~\ref{c2:ex:2} f), checked by a clean analysis
now and by \code{uart_top_tb} in \chapref{5}; and the parity and second stop bit of
Exercise~\ref{c2:ex:3} a), checked by your own testbench case.

\solution{c2:ex:1}{\code{baud_gen}}
\ans{b} \code{div = 50000000 / (16 * 115200) = 27.13}, which rounds to \textbf{27}, the value in the
protocol's table (\secref{proto:baud}). One bit is then $16 \times 27 = 432$ clocks, 8.64\,µs,
which is 115740.7~baud, 0.47\,\% fast.

The divide is by \code{16 * baud} because the receiver cannot see where a bit begins, so it samples
the line sixteen times per bit, and the transmitter and the receiver share the one time base.
\code{uart_rx} (\chapref{3}) spends the sixteen ticks: it detects the start edge on one of them,
counts to the middle of the start bit to re-check it, and then counts sixteen per bit to land each
later sample in the middle of its bit. \code{uart_tx} simply holds each bit for sixteen ticks. A
generator dividing by \code{baud} alone would give the transmitter everything it needs and leave the
receiver nothing to find the middle of a bit with.

\ans{c} The testbench passes with \code{counter = div - 1}. The counter reloads to zero on the very
edge it reaches \code{div - 1}, so with a constant \code{div} it never passes the threshold, and
``has reached'' and ``is at least'' are the same test: both fire exactly when the counter is 3.

They part company only when the counter is already past the threshold, and a \code{div} of zero is
the extreme case, since \code{div - 1} is then $-1$. With \code{>=}, every counter value is at least
$-1$, so \code{tick} is high on every clock and the counter stays at zero. With \code{=}, the
counter is never $-1$, so it never reloads: \code{tick} never rises, and the counter climbs until
its next increment leaves its declared range. Simulated with \code{div} forced to zero, the
\code{>=} version ticked on 99 of the first 100 clocks, and the \code{=} version never ticked and
stopped after 65536 clocks with

\begin{console}
ghdl:error: bound check failure at baud_gen.vhd:22
\end{console}

\noindent so \code{>=} is the one that keeps the counter inside its range. (In hardware the counter
would wrap silently instead, and still never tick.) In this design a zero cannot actually reach the
port, since \code{div} is \code{natural range 1 to 65535}: that is why \code{uart_top} substitutes
1 for a zero \code{BAUD_DIV}, and the robustness of \code{>=} is a second line of defence. It also
covers the case the range does allow, a \code{div} lowered while the counter is above the new
threshold: \code{>=} reloads on the next edge, \code{=} counts on to the top of its range first.

\ans{d} Both \code{uart_tx} and \code{uart_rx} count ticks, and both are synchronous: each looks at
\code{baud_tick} only inside \code{rising_edge(clock)}. A tick decoded combinationally from the
registered counter is still exactly one clock wide, because the counter holds \code{div - 1} for
exactly one clock. A glitch on it between edges, while the counter's bits change, is invisible to
both consumers as long as it has settled by the next edge, which is ordinary timing closure. So a
glitchy tick would not break them.

A \emph{wider} tick would. Every clock on which \code{baud_tick} is high counts as one sixteenth of
a bit, so a tick two clocks wide advances \code{uart_tx}'s \code{ticks} counter twice per period:
bits come out half as long as they should, and \code{uart_rx} samples at the wrong points. The
registered output is the design that makes ``one clock wide'' true by construction. A glitch would
matter only where \code{tick} is used as a clock or crosses into another clock domain, and it is
neither here.

\solution{c2:ex:2}{\code{uart_tx}}
\ans{b} The testbench stops at the first data bit:

\begin{console}
uart_tx_tb: data bit 0 mismatch!
\end{console}

\noindent \code{0x53} is \code{0101_0011}, so least significant first the data bits go out as
\code{1, 1, 0, 0, 1, 0, 1, 0}. Packed most significant first they go out as
\code{0, 1, 0, 1, 0, 0, 1, 1}: the first bit on the line is \code{data(7)}, a \code{0}, where the
testbench expects \code{data(0)}, a \code{1}. That is the byte shifted the other way, \code{0xCA}.

\ans{c} The testbench passes. \code{0xA5} is \code{1010_0101}. Least significant first, the eight data
levels are \code{1, 0, 1, 0, 0, 1, 0, 1}; most significant first they are
\code{1, 0, 1, 0, 0, 1, 0, 1} as well, because the byte reads the same in both directions, so the two
packings put the same ten levels on the line and a check of the levels cannot tell them apart.

For the check to be able to fail, the byte must \textbf{not be a bit-reversal palindrome}: its bit
$i$ must differ from its bit $7 - i$ for at least one $i$. A palindrome fixes the top four bits by
the bottom four, so there are $2^4 = 16$ of them, and those 16 lack the property: \code{0x00},
\code{0x18}, \code{0x24}, \code{0x3C}, \code{0x42}, \code{0x5A}, \code{0x66}, \code{0x7E},
\code{0x81}, \code{0x99}, \code{0xA5}, \code{0xBD}, \code{0xC3}, \code{0xDB}, \code{0xE7} and
\code{0xFF}.

\code{uart_rx_tb} sends \code{0x53} and \code{0x2D}, neither a palindrome, and \code{uart_top_tb}
sends \code{0x53}, so each constant is good. A loopback still cannot catch the two halves reversed
\emph{together}, however good its constant: a transmitter that sends \code{data(7)} first and a
receiver that stores the first bit it gets in bit 7 undo each other exactly, and the byte that
comes back is the byte that went in. With both halves reversed, \code{uart_top_tb} reports that it
passed. With either half reversed alone it fails with

\begin{console}
uart_top_tb: received byte mismatch, expected 0x53, got 0xCA!
\end{console}

\noindent and the receiver reversed alone fails its own testbench with the same message from
\code{uart_rx_tb}. Each unit testbench compares its module with the protocol, a fixed wire order,
rather than with the other module, so each catches its own half, and together they catch the pair
the loopback cannot.

\ans{d} \code{busy} is a level that is true for the whole frame, so it answers ``may another byte be
loaded now?'', which is the feeder's whole question; it is also half of \code{STATUS}'s TX-idle bit.
\code{done} is a one-clock event at the moment the stop bit ends, so it can count frames or raise a
transmit-complete interrupt exactly once per frame, which a level can only do with an edge detector
of its own.

A feeder that loaded on \code{done} instead of on \code{not busy} would never send anything:
\code{done} only pulses at the end of a frame, and no frame starts until something is loaded, so the
first byte waits for an event that only a sent byte can produce. Built that way, \code{uart_top_tb}
reports

\begin{console}
uart_top_tb: timed out waiting for RX valid!
\end{console}

\noindent because the byte written to \code{TX_DATA} never leaves the FIFO. Even with a first load
forced some other way, a one-clock pulse is the wrong thing to gate on: miss it, because the FIFO
was empty in that one clock and a byte arrived a cycle later, and that byte waits for a frame that
will never end.

\ans{e} \code{frame <= '1' & data & '0'} puts the leftmost item in the highest bit, so

\begin{narrowtable}{@{}lcccccccccc@{}}
  \thead{\code{frame} index} & 9 & 8 & 7 & 6 & 5 & 4 & 3 & 2 & 1 & 0 \\
  \midrule
  \thead{Bit} & stop & \code{d7} & \code{d6} & \code{d5} & \code{d4} & \code{d3} & \code{d2} &
    \code{d1} & \code{d0} & start \\
  \thead{Value} & \code{1} & & & & & & & & & \code{0} \\
\end{narrowtable}

\noindent Read from index 0 upward, that is the start bit, then \code{d0}, \code{d1} and so on to
\code{d7}, then the stop bit: exactly the wire order of an 8N1 frame. The caller may change
\code{data} on the very next cycle because \code{frame} is a register of the transmitter's own,
loaded once, in \code{STATE_IDLE}, on the edge that sees \code{start}. From then on every bit on the
line comes from \code{frame}, and \code{data} is not read again until the next \code{start} in the
next idle period. In \code{uart_top} that is what lets the FIFO advance on the same edge the
transmitter loads.

\solution{c2:ex:3}{Parity and a second stop bit}
\ans{b} The stop field, which is now one or two bits long. (With the parity bit of part a), the
parity field is present or absent too, so the frame is 10, 11 or 12 bits; the start bit and the
eight data bits never change.) At a fixed baud rate the frame time is the number of bits over the
baud rate, so a second stop bit makes every frame one bit period longer, 11 instead of 10: at
115200~baud that is 95.5\,µs instead of 86.8\,µs, and the throughput falls by a tenth,
from 11520 to 10473 bytes per second.

\ans{c} \code{uart_tx_tb} binds positionally and names eight ports, so it binds the eight existing
ones and leaves the three new inputs, declared after them, at their defaults of \code{'0'}. With
\code{parity_en} and \code{two_stop} low the frame is the old 8N1 frame, bit for bit, and the
testbench passes unchanged. (Declared anywhere else, or without a default, a new input shifts or
breaks the positional binding, and the testbench no longer analyzes against the module.)

To pin down even parity, drive \code{parity_en = '1'} and \code{parity_odd = '0'}, send \code{0x53},
and after the eight data bits sample one more bit centre: expect \code{'0'}, because \code{0x53}
has four 1s, already an even number. Then sample one bit later and expect the stop bit's
\code{'1'}. The choice of a byte whose parity bit is \code{'0'} is deliberate: a transmitter that
forgot to insert the parity bit would put its stop bit, a \code{'1'}, in that position, and one that
computed odd parity would send a \code{'1'} too, so both fail. Add a second byte of odd weight, such
as \code{0x52} with three 1s, and expect a \code{'1'}, so that a parity bit stuck at \code{'0'}
fails as well.

% Section B.3: the answers to the parts of lectures/L03/appendix/c_exercises.md that are not code.
% Every testbench message quoted was produced by a design built to Chapters 1 to 5 and broken in
% exactly the way the part describes.

\section{Chapter 3: The Synchronizer and the Receiver}
\label{sol:c3}
The a) parts of both exercises are code: \code{sync}, checked by \code{sync_tb}, and
\code{uart_rx}, checked by \code{uart_rx_tb}.

\solution{c3:ex:1}{\code{sync}}
\ans{b} The first two-edge check fails and the run stops:

\begin{console}
sync: output changed after only one clock edge!
\end{console}

\noindent It says exactly what was deleted: the input reached the output one edge after it changed,
so there is one register between them where there should be two.

Every other simulation in the course would pass, because \textbf{nothing in a simulator is ever
metastable}. A one-flop synchronizer is, functionally, a correct synchronizer with one cycle less
latency, and no other testbench cares about one cycle: \code{uart_rx_tb} drives \code{rx_s2}
directly and never instantiates \code{sync}, and in \code{uart_top_tb} the looped-back line changes
in step with the same clock, so it never violates anything. The bug is purely electrical: on the
bench, the single flop samples a genuinely asynchronous line, sometimes inside its setup and hold
window, and its output occasionally hangs between levels for long enough to reach the receiver's
logic. Only a testbench that checks the latency itself can see the missing flop, which is why
\code{sync_tb} checks it.

\ans{c} Cases 0 to 4 still pass: the reset is held low across two clock edges there, so a synchronous
reset has time to clear the chain. Case 5 asserts \code{reset_s2_n} a quarter of a period after an
edge with \code{1011} standing on the wide output and looks one nanosecond later:

\begin{console}
sync: reset_s2_n did not clear the output asynchronously!
\end{console}

\noindent The assertion has to be asynchronous so that the chain clears the moment reset is asserted,
whether or not a clock edge follows: a design is reset when its clock has not started yet, or has
stopped, and every other block in \file{hw/} clears the same instant, so a block that waited for an
edge would be the one block out of step. The release, on the other hand, is not \code{sync}'s
problem. \code{reset_s2_n} comes from \code{reset_sync}, which already lines its rising edge up with
the clock, two flip-flops after \code{reset_n} rises. By the time the release reaches \code{sync}
it is an ordinary synchronous signal, so testing it in the sensitivity list is safe.

\ans{d} A module that needs a signal as an input cannot also be where that signal comes from.
\code{sync} resets its flops from \code{reset_s2_n}; to produce \code{reset_s2_n} it would have to
be reset by its own output, which is undefined until the module has already run, and the raw
\code{reset_n} would have to go through its data input instead, where the assertion would take two
clock edges to arrive and would not arrive at all without a clock. \code{reset_sync} is built the
other way round: the raw \code{reset_n} is the asynchronous clear of both its flops, and a constant
\code{'1'} is clocked through them, so the assertion needs no clock and the release is synchronous.
The two modules look alike, two flops in series, and are wired in opposite directions, so the course
keeps both.

\ans{e} Each bit of a wide \code{sync} is its own two-flop chain. If two input bits change on the same
edge of their own domain and the change lands inside the first flops' setup and hold window, each
first flop may go metastable and each resolves independently, one to the new value and one to the
old. The one that resolved to the new value appears at the output two edges later, the other three
edges later, and for one clock \code{sync_out} shows one bit's new value beside the other's old one.

That is not a bug in \code{sync}: it promises that every bit it delivers is a clean level, and that
each input change arrives within two or three edges. It does not promise that bits which changed
together arrive together, and nothing built from independent flops can. For eight unrelated pins it
does not matter, since each pin is read on its own. For a counter it is fatal: going from
\code{0111} to \code{1000}, every bit changes, and a mixed sample can be any of sixteen values,
including \code{1111} and \code{0000}, which the counter never held. A multi-bit value crosses
domains as a Gray code, where one bit changes per step, or through a handshake or an asynchronous
FIFO, never through a wide synchronizer.

\solution{c3:ex:2}{\code{uart_rx}}
\ans{b} \textbf{The stop bit} drifts furthest. The receiver resynchronizes only on the start edge; every
later sample is timed by counting its own ticks from there, so a rate error accumulates with the
time since that edge. The first data bit is sampled about 1.6 bit periods after the edge, the stop
bit about 9.7, so the stop bit carries six times the drift.

With the transmitter faster by a fraction $f$, its bit period is $1/(1+f)$ of the receiver's. Its
stop bit, the tenth bit of the frame, ends ten of its own bit periods, 160 of its ticks, after the
edge. The stop-bit sample has to land before that.

\begin{itemize}
  \item \textbf{An ideal tick-8 sampler} samples the stop bit at the centre of its tenth bit, 9.5 bit
    periods, 152 of its ticks, after the edge. It needs $152 < 160/(1+f)$, so $f < 160/152 - 1 =
    5.3\,\%$. Counting the up to one tick a 16-times sampler inherently takes to notice the edge,
    $f < 160/153 - 1 = 4.6\,\%$.
  \item \textbf{This receiver} samples the stop bit 155 ticks after the tick that noticed the edge
    (9 to the start re-check, 17 to the first data bit, 16 for each of the other seven, 17 to the
    stop bit), and the edge is noticed up to one tick after it happens, so up to 156 ticks after the
    edge. It needs $156 < 160/(1+f)$, so $f < 160/156 - 1 = 2.6\,\%$.
\end{itemize}

\noindent So the transmitter can run about 5\,\% fast against the ideal sampler, and only about
2.5\,\% fast against this one, which spends half the budget on its own late sampling phase before
any mismatch starts. The budget is shared by both ends: the DE0-CV's divider alone already runs
0.47\,\% fast at 115200. The same late phase buys margin the other way: a transmitter up to
$155/144 - 1 = 7.6\,\%$ \emph{slower} still has its stop bit under the sample.

\ans{c} With the re-check removed, a single tick of low on an idle line starts a frame. The receiver
then samples a line that has already gone back to idle, reads eight \code{1}s and a high stop bit,
and delivers a perfectly well-formed \code{0xFF} with \code{valid} and no framing error. Simulated
with a one-tick glitch, the receiver without the re-check delivered one byte, \code{0xFF}, and no
framing error; with the re-check it delivered nothing. Note that \code{uart_rx_tb} still passes
without the re-check: none of its three cases contains a glitch, so the bench cannot see the
difference.

It is more than theoretical because real lines are noisy, and the moments they are noisiest are the
ones nobody is watching: a cable plugged in or pulled out, the far end powering up or down, a long
unshielded run beside a motor or a switching supply, a receive pin left floating with no pull-up.
Each of those can put a short low pulse on an idle line, and without the re-check each pulse is a
phantom \code{0xFF} in the RX FIFO.

\ans{d} \code{0x53} is \code{0101_0011}, so least significant first the bits arrive in this order:

\begin{narrowtable}{@{}lcccccccc@{}}
  \thead{Arrival} & 1st & 2nd & 3rd & 4th & 5th & 6th & 7th & 8th \\
  \midrule
  \thead{Value} & \code{1} & \code{1} & \code{0} & \code{0} & \code{1} & \code{0} & \code{1} &
    \code{0} \\
  \thead{\code{bit_idx}, index stored} & 0 & 1 & 2 & 3 & 4 & 5 & 6 & 7 \\
\end{narrowtable}

\noindent so \code{frame(7 downto 0)} holds \code{0}, \code{1}, \code{0}, \code{1}, \code{0},
\code{0}, \code{1}, \code{1} from bit 7 down, which is \code{0101_0011}, \code{0x53}. Counting down
from 7 instead stores the first bit in bit 7, so \code{frame} holds \code{1100_1010}: \code{data_out}
reads \code{0xCA}, the \textbf{bit reversal} of the byte sent, and \code{uart_rx_tb} says so:

\begin{console}
uart_rx_tb: received byte mismatch, expected 0x53, got 0xCA!
\end{console}

% Section B.4: the answers to the parts of lectures/L04/appendix/b_exercises.md that are not code.
% Every testbench message quoted was produced by a design built to Chapters 1 to 5 and broken in
% exactly the way the part describes.

\section{Chapter 4: The FIFO and the Receive Path}
\label{sol:c4}
Exercise~\ref{c4:ex:1} a) is code, \code{fifo}, checked by \code{fifo_tb}; so is all of
Exercise~\ref{c4:ex:2}, the receive path in \code{uart_top}, which is checked by a clean analysis
now and by \code{uart_top_tb} in \chapref{5}.

\solution{c4:ex:1}{\code{fifo}}
\ans{b} After four pushes \code{head} has wrapped round to 0, equal to \code{tail}, and
\code{count} is 4. The unguarded fifth push stores \code{0x99} in entry 0, over \code{0x11}, the
oldest byte, which has not been read; advances \code{head} to 1, one past \code{tail}; and raises
\code{count} to 5. The FIFO now claims five entries in a four-entry ring, and its front byte is
gone. What the run reports depends on the declaration of \code{count}:
\begin{itemize}
  \item As \code{natural range 0 to DEPTH}, the assignment of 5 is out of range, and GHDL stops on
    the fifth push's edge, before any assertion runs, with a
    \code{ghdl:error: bound check failure at fifo.vhd:}\emph{n}, where \emph{n} is the line that
    increments \code{count}.
  \item As an unconstrained \code{natural}, \code{count} becomes 5, and \code{full}, which is
    \code{count = DEPTH}, drops. The testbench's own check fires:

\begin{console}
fifo_tb: still full after a dropped write!
\end{console}
\end{itemize}
Both are failures, and only the second is the testbench talking; the first is the language refusing
a value the declaration said could not happen, which is what a constrained type is for.

\ans{c} Because the register bus has no read strobe. A read in this design is not an event: the
bridge sets \code{reg_addr}, the bank's read mux presents \code{reg_rdata} combinationally, and the
bridge latches the word once, one cycle after the command byte. The bank never learns that a read
has happened, let alone when it finished. A FIFO whose only read operation was ``read and pop''
would force the bank to invent a pop out of that non-event, and then decide what it means when the
transaction is abandoned half-way. Look-then-advance lets \code{RX_DATA} be wired straight to
\code{rdata}, a pure read with no side effect, like \code{STATUS}, and puts the advance on \code{rd},
where the separate \code{RX_POP} \emph{write} drives it. Writes are exactly what the bridge already
commits once, on completion, and never on an abort, so the pop inherits that for free.

\ans{d} \textbf{Eight.} Eight back-to-back frames fill the depth-8 RX FIFO with nothing read; the
ninth is the first that can be lost. The receiver holds no finished byte of its own: it pushes each
byte on the one-clock \code{valid} at its stop-bit sample and keeps only the frame it is still
assembling. So while the ninth frame is on the wire, software has until that frame's stop-bit sample
to pop one byte, about one frame time, 86.8\,µs at 115200~baud; if the FIFO is still full
then, the ninth byte is dropped by the FIFO's \code{not full} guard, and every frame after it too.

\textbf{No flag tells software, in this build.} \code{STATUS} bit 1 stays set throughout, which says
only that there is something to read, and \code{ER_OVERRUN} has no producer, so it reads 0 (see
\secref{proto:reserved}). The information exists: \code{uart_regs} exports \code{rx_full}, and a
\code{valid} arriving while it is high is precisely a lost byte. Exercise~\ref{c4:ex:3} is about
where to latch it.

\solution{c4:ex:3}{Overrun belongs to the register bank}
\ans{a} An overrun is a byte arriving when there is nowhere to put it, and whether there is anywhere
to put it depends on whether the previous bytes have been read. \code{uart_rx} does not know that. It
emits a one-clock \code{valid} and forgets the byte; it has no view of the FIFO behind it or of the
software behind that. The one piece of information it lacks is \textbf{the consumer's state}, which
here is the RX FIFO's \code{full} flag.

\ans{b} Add a third latch beside the framing one in \code{uart_regs}: in the clocked process, after
the write decode, set \code{err_flags(ER_OVERRUN)} when \code{rx_push} and \code{rx_full_s} are both
high, since that is a push the FIFO is about to drop. It lives in \code{ERROR_FLAGS} bit 2, at the
position the protocol reserves for it, latched until software writes \code{0x0}, and \code{STATUS}
bit 2 shows it with no further change, because that bit is already the OR of the three flags. The
driver's existing \code{errorFlags()} and \code{clearErrors()} then reach it unchanged, which is what
reserving the position bought.

\ans{c} A \textbf{framing error} is a wire problem: the stop bit was low when it was sampled. The
causes are on the line or in its configuration: the two ends at different baud rates, a divider
written wrongly, noise on the stop bit, a sender using a different frame format such as a parity bit
where the receiver expects the stop bit, or a break, the line held low because a cable came loose
or the transmitter lost power. An \textbf{overrun} is a software timing problem: every frame on the
wire was perfect, and the driver simply did not read them as fast as they came, because it was busy
elsewhere, blocked in a delay, or polling too slowly. A framing error says do not trust this byte; an
overrun says a good byte was thrown away.

% Section B.5: the answers to the parts of lectures/L05/appendix/b_exercises.md that are not code.
% Every testbench message quoted was produced by a design built to Chapters 1 to 5 and broken in
% exactly the way the part describes.

\section{Chapter 5: The Register Bank}
\label{sol:c5}
The a) parts of both exercises are code: \code{uart_regs}, checked by \code{uart_regs_tb}, and the
completed \code{uart_top}, checked by \code{uart_top_tb}.

\solution{c5:ex:1}{\code{uart_regs}}
\ans{b} Advancing the FIFO ``whenever that index is read'' means decoding a pop from the address
alone, since the bus has no read strobe: \code{rx_pop} is high whenever \code{reg_idx} is
\code{REG_RX_DATA}. The run then fails one check earlier than you might expect, in Case~5 rather
than Case~6:

\begin{console}
uart_regs_tb: RX_DATA mismatch, got 0x00!
\end{console}

\noindent \code{bus_set_addr} puts the address on the bus, waits for a clock edge and then samples.
That edge pops the FIFO, so by the time the bench reads \code{RX_DATA} the byte it pushed is gone
and the front shows an entry that was never written. The ``a bare \code{RX_DATA} read must not pop''
check in Case~6 would fail too, but the run never gets there.

The system testbench, by contrast, \emph{passes} with the change. The bridge latches
\code{reg_rdata} on the same edge as the first pop, so it captures the byte before it disappears,
and \code{uart_top_tb} reads \code{RX_DATA} once and never looks again. What it does not see is the
rest: the bridge holds \code{reg_addr} at the last command's index until the next command, so the
pop repeats on every clock and an eight-entry FIFO is empty 160\,ns after the command byte.

That is the general problem, and the abort rule makes it plain. A read with a side effect has to
decide \emph{when} the side effect happens. If the pop happens when the value is latched, after the
command byte, then an aborted read, \code{SS} rising before the fifth byte, has discarded a byte the
master never received, and the protocol promises that an abort has no side effects. If the pop waits
for the fifth byte, it is a commit on completion, the discipline a write has; but the bridge commits
only writes, with \code{reg_write}, and has no ``read completed'' strobe to give the bank, so both
the bridge and the bank would have to change. The pure read plus a separate \code{RX_POP} needs
none of that: reading has no side effect to abort, and the pop is a write, which the bridge already
commits exactly once, on completion, and never on an abort.

\ans{c} \code{STATUS} bit 0, TX ready, is \code{not tx_full_s}, from the TX FIFO. Bit 1, RX valid, is
\code{not rx_empty_s}, from the RX FIFO. Bit 2, Error, is the OR of the three \code{err_flags}
latches. Bit 3, TX idle, is \code{tx_empty_s and not tx_busy}, from the TX FIFO and the
transmitter. Bits 31 to 4 are zero.

None of them can be a stored register the datapath writes, because each is a copy of state that
already lives somewhere else: the FIFOs' counts, the error latches, the transmitter's state machine.
A stored copy would have to be updated by every event that changes its source, on the same edge,
by logic that duplicates the FIFO's own; any cycle in which it lagged would be a cycle in which the
driver could read ``TX ready'' for a FIFO that had just filled and push a byte into it that is then
dropped. Computed combinationally, each bit is correct by construction, and the bridge's latch-once
read takes a coherent snapshot of all four at once.

\ans{d} It relies on the bridge raising \code{reg_write} only for a \textbf{completed, non-aborted
write}: one clock, once, when the fifth byte arrives with \code{SS} still low, with \code{reg_addr}
and \code{reg_wdata} already holding the transaction's index and all four data bytes. The provided
\code{spi_reg_bridge} does exactly that: it counts the bytes, resets the count whenever \code{SS}
rises, and asserts \code{reg_write} only on byte four of a write command, which
\code{spi_reg_bridge_tb} checks, abort included.

If a half-finished transaction could reach the bank as a \code{reg_write}, its data word would be
partly stale, the bytes left over from an earlier transaction, and the bank would act on it: a
\code{BAUD_DIV} of garbage, and the line at a rate nobody chose; a \code{TX_DATA} push of a byte the
master never finished sending, which then goes out on \code{tx}; an \code{RX_POP} that discards an
unread byte; an \code{ERROR_FLAGS} write that clears an error nobody saw. Every one of those is a
single-cycle action with nothing to undo it, which is why the guarantee has to hold before the
strobe reaches the bank.

\solution{c5:ex:2}{\code{uart_top}}
\ans{b} A release that lands near a clock edge would leave some flip-flops out of reset on that edge
and others on the next, with any of them possibly metastable, and the peripheral would start
half-reset. Here that could mean a FIFO whose \code{head} leaves reset a cycle before its
\code{count}, a receiver whose \code{rx_prev} leaves reset after its state and sees a falling edge
that never happened, or the transmitter's state machine out of step with its \code{tx} register:
states the logic was never designed to leave. Asserting asynchronously is nonetheless right, because
entering reset has no such hazard. Every flop goes to its reset value immediately, together, and
without needing a clock, so the design is safe the instant \code{reset_n} falls, even if the clock
has stopped, and \code{tx} goes to its idle high at once. Only the release has to line up with the
clock, and \code{reset_sync} lines it up.

\ans{c} The port that settles it is \code{sync}'s \code{reset_s2_n} input. To make
\code{reset_s2_n} with a \code{sync} instance, you would have to connect the instance's own output,
\code{reset_s2_n}, to its own reset input: the synchronizer would be reset by the signal it has not
produced yet, which is undefined until it has been clocked, and it would never assert while it was
holding itself in reset. The raw \code{reset_n} would have to go in through \code{async_in}, so even
the assertion would wait two clock edges and would never come at all without a clock. The circle
cannot be closed; \code{reset_sync} breaks it by taking the raw reset on its flops' asynchronous
clear.

\ans{d} \code{tx_load} is \code{(not tx_empty) and (not tx_busy)}, and \code{tx_pop} is the same
signal.
\begin{itemize}
  \item \textbf{Never zero.} Whenever a byte is queued and the transmitter is idle, \code{tx_load} is
    already high, combinationally, so the transmitter loads on the next edge.
  \item \textbf{Exactly one.} On that edge, \code{uart_tx} sees \code{start} in \code{STATE_IDLE},
    latches the front byte and moves to \code{STATE_SEND}, and the FIFO pops the same byte. Because
    \code{busy} is decoded combinationally from the state register, it is high from that edge on,
    so \code{tx_load} falls in the very next cycle. It cannot reload while the transmitter is busy,
    and the FIFO cannot pop a second byte for the same frame.
  \item \textbf{Held low} for the rest of the frame by \code{tx_busy}, which stays high from the
    load edge until the edge on which the stop bit's sixteenth tick returns the state machine to
    \code{STATE_IDLE}. In that first idle cycle, a queued byte raises \code{tx_load} again.
\end{itemize}
Simulated with three bytes queued at once, the feeder popped exactly three times, once per frame,
each pop on the first idle edge after the previous frame's \code{done}. The argument depends on
\code{busy} rising on the load edge itself. With \code{busy} registered one cycle late, as a second
flip-flop, \code{uart_tx_tb} and \code{uart_top_tb} both still pass, since each sends one byte and a
second pop from an empty FIFO is ignored; but with three bytes queued the feeder popped twice on
consecutive edges, the transmitter took the first and ignored the second, and the third byte never
arrived.

\ans{e} One byte, around the loop:
\begin{enumerate}
  \item The testbench plays a five-byte write of \code{TX_DATA} on \code{sclk}, \code{mosi} and
    \code{ss}. \code{spi_slave} synchronizes the three lines and shifts each byte in;
    \code{spi_reg_bridge} assembles the command and the four data bytes and, on the fifth, raises
    \code{reg_write} with \code{reg_addr} 3 and the byte in \code{reg_wdata(7 downto 0)}.
  \item \code{uart_regs} decodes that as \code{tx_push}, and its TX \code{fifo} stores the byte.
  \item The \textbf{feeder} in \code{uart_top} sees the FIFO not empty and the transmitter not busy,
    and raises \code{tx_load}: \code{uart_tx} latches the byte into its frame, and the TX FIFO pops.
  \item \code{uart_tx} drives the frame out on \code{tx}, one bit per sixteen \code{baud_gen}
    ticks.
  \item The testbench's loopback, \code{rx <= tx}, returns it on the \code{rx} pin.
  \item \code{sync} carries \code{rx} into the clock domain as \code{rx_s2}.
  \item \code{uart_rx}, on the same ticks, finds the start edge, samples the eight bits and the stop
    bit, and pulses \code{valid}, which is \code{rx_push}.
  \item \code{uart_regs} pushes the byte into its RX \code{fifo}, and \code{STATUS} bit 1 rises.
  \item The testbench polls \code{STATUS} over SPI until it sees bit 1, then reads \code{RX_DATA}:
    \code{spi_reg_bridge} latches \code{reg_rdata}, which the bank's read mux fills from the RX
    FIFO's front, and \code{spi_slave} shifts it out on \code{miso}, most significant byte first.
\end{enumerate}
The system testbench stops there, with the byte still in the RX FIFO: it never writes \code{RX_POP}.

% Section B.9, from lectures/L09/appendix/c_exercises.md: every part that is not code. The outcomes
% quoted were measured on a reference build of the whole stack made to the chapters'
% specifications: the host suites over the mocked register file, and the avr-gcc 7.3.0 image and
% its disassembly.

\section{Chapter 9: The Real Transport, and Both Toolchains}
\label{sol:c9}
The code in this chapter is \code{AvrSpi} (Exercise~\ref{c9:ex:1} a), checked by the provided
suite \file{fw/test/transport/avr_spi_test.cpp} over the mocked register file, and the bring-up
\code{main} and the logging (Exercise~\ref{c9:ex:2} a and Exercise~\ref{c9:ex:3} a), which the
bench checks. Everything else is answered here.

% ------------------------------------------------------------------------------------------
\solution{c9:ex:1}{\code{AvrSpi}}
\ans{b} \code{SPDR} holds the incoming byte only once all eight bits have been shifted, eight
\code{SCK} periods after the write, which at f\textsubscript{osc}/16 is 128 CPU cycles, and
\code{SPIF} is the only thing that says that moment has come: read earlier and you get a stale
byte, write the next byte earlier and you collide with the transfer still in progress.

The spin is cheaper than it looks. Compiled, the whole of \code{transfer()} is six instructions,
and the wait is a four-cycle loop that goes round about thirty times per byte:

\begin{console}
out  0x2e, r22      ; SPDR = byte: the transfer starts
in   r0, 0x2d       ; read SPSR
sbrs r0, 7          ; SPIF set? skip the jump
rjmp .-6            ; not yet: read SPSR again
in   r24, 0x2e      ; return SPDR
ret
\end{console}

\ans{c} If \code{SS} (PB2) is an input and something drives it low while \code{MSTR} is set, the
SPI peripheral concludes that another master has taken the bus and gets out of its way:
\textbf{it clears \code{MSTR} in \code{SPCR}}, demoting itself to a slave, and sets \code{SPIF},
raising the SPI interrupt if one is enabled. The ATmega328P has no separate mode-fault flag;
\code{MSTR} reading back as 0 is the only record of what happened.

For \code{AvrSpi} that is two failures in a row. The \code{transfer()} in progress sees \code{SPIF}
and returns whatever is in \code{SPDR}. Every \code{transfer()} after it writes \code{SPDR} as a
slave, where nothing is driving \code{SCK}, so \code{SPIF} never sets again and the poll spins for
ever: the driver hangs in the first transaction after the glitch, and nothing sets \code{MSTR} back.
Making PB2 an output removes the possibility, and costs nothing, because the transport drives that
pin as the chip select anyway.

\ans{d} It returns the byte the \emph{previous} transfer received, still sitting in \code{SPDR}'s
receive buffer, or whatever was there after reset if there was no previous transfer. The new byte
does not exist yet: it arrives eight \code{SCK} periods later, 8\,\textmu s or 128 CPU cycles after
the write, and the read happens one cycle after it. Every reply is therefore shifted by one byte,
so a \code{readReg} assembles the command byte's reply and three data bytes instead of four.

On the part it is worse than a shift. The next \code{transfer()} writes \code{SPDR} a few cycles
later, while the byte before it is still shifting, which the datasheet calls a write collision:
\code{WCOL} is set and the write is \emph{ignored}, so most of the bytes of a transaction never
leave the master at all, and the five-byte transaction is garbled in both directions.

The mock makes the same point more bluntly. It completes a transfer only when the driver reads
\code{SPSR}, so with the poll gone no transfer ever completes and every read of \code{SPDR} returns
the value left from reset. Three cases go red, and only these three:

\begin{console}
Test case AvrSpi.TransferExchangesByte failed: EXPECT_EQ(rx, miso) failed: 0 != 205
Test case AvrSpi.TransferIsFullDuplexInOrder failed: EXPECT_EQ(spi.transfer(mosi[i]), miso[i])
  failed: 0 != 170
Test case AvrSpi.DrivesTheDriverEndToEnd failed: EXPECT_EQ(uart.status(), data) failed:
  0 != 305419896
19 out of 22 test cases succeeded!
\end{console}

\noindent Every \code{Uart} case stays green, because those run over the scripted
\code{transport::Stub}, not over \code{AvrSpi}: the bug is below the seam, and so is the only suite
that can see it.

% ------------------------------------------------------------------------------------------
\solution{c9:ex:2}{The freestanding bring-up \code{main}}
\ans{b} Not the provided demo itself: \file{source/main.cpp} was written in the AVR-portable
subset, and it compiles and links unchanged with the \file{fw/avr} flags, into 1256 bytes of flash
and 34 of RAM. It needs none of the idioms a freestanding build removes: its objects are automatic,
with fixed arrays inside; the driver reports failure with a \code{bool}; and it prints nothing.

It is the idioms of ordinary host C++ that fail, and each has a replacement on the AVR:

\begin{booktable}{@{}p{10.6em}p{11.4em}L@{}}
  \thead{Host idiom} & \thead{In the AVR build} & \thead{What replaces it} \\
  \midrule
  \code{#include <cstdint>}, \code{std::uint8_t} & Does not compile: avr-libc has no C++ library,
    so \code{fatal error: cstdint: No such file or directory} & \code{<stdint.h>} and the bare
    \code{uint8_t} \\
  \code{new}, \code{std::make_unique}, \code{std::vector} & No \code{operator new}, and no standard
    containers & Static and automatic storage, sized at compile time \\
  \code{throw}, \code{try} & Rejected under \code{-fno-exceptions} & Return values, as
    \code{write()} returning \code{false} \\
  \code{dynamic_cast}, \code{typeid} & Rejected under \code{-fno-rtti} & Virtual dispatch through
    an interface \\
  \code{std::cout}, \code{printf} to a console & No iostreams, and no console & Bytes written to the
    ATmega's own USART, as in Exercise~\ref{c9:ex:3} \\
\end{booktable}

\ans{c} With the reference firmware, the linker reports one symbol, three times:

\begin{console}
In function `app::EchoNode::~EchoNode()':
  undefined reference to `operator delete(void*, unsigned int)'
In function `driver::transport::AvrSpi::~AvrSpi()':
  undefined reference to `operator delete(void*, unsigned int)'
In function `driver::uart::Uart::~Uart()':
  undefined reference to `operator delete(void*, unsigned int)'
\end{console}

\noindent The construct is the \textbf{virtual destructor}. Every class that has one gets a second,
``deleting'' destructor, which runs the destructor and then calls \code{operator delete}, so that
\code{delete} through a base pointer frees the right object; the vtable points at it, so it is
emitted, and names \code{operator delete}, even though nothing in the program ever deletes anything.
From C++14 it calls the sized form, \code{operator delete(void*, size_t)}, and \code{size_t} is an
\code{unsigned int} on the AVR. That is why the linker names the three concrete classes: each
inherits a virtual destructor from its \code{Interface}.

The other symbols \file{env.cpp} defines turn up with other flags and other code, which is why it
defines them all. At \code{-O0} the same link also misses \code{__cxa_pure_virtual}, eight times:
the \code{Interface} classes' own vtables name it for every pure virtual, and only the optimizer
removes those vtables, because every \code{Interface} is constructed as the base of a concrete class
that overwrites its vtable pointer at once. Without \code{-fno-threadsafe-statics}, a
function-local \code{static} with a constructor, such as a \code{static driver::uart::Uart
uart\{spi\};} inside \code{main()}, adds \code{__cxa_guard_acquire} and \code{__cxa_guard_release}.
And under \code{-fno-sized-deallocation} the missing symbol is the unsized
\code{operator delete(void*)} instead.

The host build does not need the file because g++ links libstdc++, which defines every one of these
symbols, its static guards thread-safe. \file{env.cpp} lives in \file{avr/}, where the host Makefile
never looks, precisely so that its single-threaded guards can never replace those.

% ------------------------------------------------------------------------------------------
\solution{c9:ex:3}{Logging and timing}
\ans{b} Five bytes of eight bits is \textbf{40 \code{SCK} cycles}, which at 1\,MHz is \textbf{40\,\textmu
s} of clocking. The analyzer shows more: five bursts of eight clocks, each 8\,\textmu s long, with a
gap between consecutive bursts, and \code{SS} falling a little before the first and rising a little
after the last.

The gaps are software. \code{SCK} runs only while a byte is shifting, and between two bytes the CPU
has to notice \code{SPIF}, return from \code{transfer()}, go round \code{readReg}'s loop and fold the
byte into its 32-bit result, make the next virtual call through the transport interface, and write
\code{SPDR} again. Counted from the reference build's disassembly, that path is about 30 CPU cycles,
roughly 2\,\textmu s at 16\,MHz, so a register read takes about $5 \times 8 + 4 \times 2 = 48$\,\textmu
s between the first clock and the last, and about 50\,\textmu s from \code{SS} falling to \code{SS}
rising, once the \code{begin()} and \code{end()} calls on either side are counted. Your gaps depend on
your code and your compiler, but not by more than a microsecond or two; what does not change is that
the bus clocks for about four fifths of a transaction and waits for the CPU for the rest.

With the logging from part a) in the loop, the gaps \emph{between} transactions are a different
matter: a twenty-character line at 115200~baud takes 1.7\,ms to leave the USART, more than thirty
register reads.

\ans{c} The SPI is \textbf{synchronous}. The master's hardware makes \code{SCK} itself, by dividing
the CPU clock by the prescaler chosen in \code{SPCR} and \code{SPSR}, and sends it on its own wire;
the slave samples on the edges it receives. The code chooses only a ratio, f\textsubscript{osc}/16,
never a frequency, and the exact rate does not matter to either end so long as it stays within what
the slave and the wiring can take. No number in the transport depends on the clock, so it needs no
\code{F_CPU}.

The USART is \textbf{asynchronous}. There is no clock wire: the PC's end has a clock of its own and
agrees with the ATmega only on a number, the baud rate, which the ATmega has to produce from its own
clock by computing a divider, \code{UBRR0 = F_CPU / (16 * baud) - 1}, rounded to an integer. The code
can only compute that from the clock frequency, so \code{F_CPU} has to be the real one: with a wrong
\code{F_CPU} every character leaves at the wrong rate and arrives as garbage.

The rounding matters at the log's end of the line too. At 16\,MHz, 115200~baud needs a divider of
7.68, which rounds to 8 and gives 111111~baud, 3.5\,\% slow, at the edge of what a frame tolerates.
Setting \code{U2X0}, which halves the divisor's factor from 16 to 8, gives a divider of 16 and
117647~baud, 2.1\,\% fast; 38400~baud (divider 25) and 9600~baud (divider 103) are both within
0.2\,\%.

% ------------------------------------------------------------------------------------------
\solution{c9:ex:4}{Your own \code{uart_top} on the board}
\ans{a} The \textbf{Fitter} report's resource usage summary gives the size: the logic utilization
in ALMs, out of the 18\,480 of the DE0-CV's Cyclone~V (a 5CEBA4F23C7), the registers, and the pins.
This peripheral is a small fraction of the device, well under a few per cent; the exact figure
depends on your code and on the Quartus version, and is worth recording mainly so that you notice
if a later change makes it jump.

The \textbf{Timing Analyzer} gives the other number. Read the worst-case \textbf{setup slack} and
\textbf{hold slack} for the 50\,MHz clock across all the timing corners Quartus analyzes, slow and
fast silicon at the hot and the cold end of the temperature range; the multicorner summary gathers
the worst of them in one table. \textbf{Positive slack on every path, in every corner, means the
design met timing}; the restricted \code{Fmax} for the clock, which should come out comfortably
above 50\,MHz for a design this shallow, says the same thing another way. Check first that the
clock is actually constrained, a 20\,ns period on the 50\,MHz pin in the project's timing
constraints: an unconstrained clock is analyzed against nothing, and its slack means nothing. The
asynchronous inputs, \code{rx} and the three SPI inputs, are expected to show up as unconstrained,
because they enter through synchronizers and have no timing relationship to the clock to check.

\emph{The report} is how you know, and the design working is not, because the analysis covers the
slowest silicon, the lowest supply voltage and the highest temperature the part is specified for,
and the board on your desk is one sample, at room temperature, today. A design with negative slack
can work on the first try and fail warm, or on the next board; and nothing a working design does
tells you how much margin it has left.

\ans{c} The pin loopback proves the \textbf{physical path}, which the simulation never had in it:
that the FPGA is configured at all, that the wrapper's pin assignment puts \code{rx} and \code{tx} on
the header pins the wires go to, and that the adapter really is at 3.3\,V logic. It cannot tell you that the terminal's rate is
right: the adapter sends and receives at the same setting, so a loopback agrees with itself at any
rate, and the rate is first on trial when the peripheral talks to the terminal in \chapref{10}. No simulation has a pin, a voltage, a USB adapter
or a terminal in it.

\code{uart_top_tb} proves the \textbf{peripheral}, which the pin loopback never had in it: that a
byte written to \code{TX_DATA} over SPI leaves \code{uart_tx} correctly framed at the rate
\code{BAUD_DIV} sets, and comes back through \code{sync}, \code{uart_rx} and the RX FIFO to be read
from \code{RX_DATA}. With \code{rx} wired straight to \code{tx} in the wrapper, none of your logic is
in the path, so a character echoing in the terminal says nothing about it. Neither of the two says
anything about timing, which is what part a) is for.

% Section B.10, from lectures/L10/appendix/b_exercises.md and c_bringup.md: every part that is not
% code. The EchoNode outcomes quoted were measured by running the provided test, and an empty-queue
% case, against a reference EchoNode and against deliberately broken ones.

\section{Chapter 10: Integration and Bring-Up}
\label{sol:c10}
\code{app::EchoNode} and its tests (Exercise~\ref{c10:ex:b1} a, b and d) are code, checked by
\file{fw/test/app/echo_node_test.cpp} and the case you add to it. Assembling the bench, programming
the two chips and climbing the ladder (Exercises~\ref{c10:ex:c1} to~\ref{c10:ex:c3}) are procedures
whose check is the bench itself: each rung either behaves as its description says, or names the
layer at fault. What is answered here is the reasoning in Exercise~\ref{c10:ex:b1} and the closing
trace.

% ------------------------------------------------------------------------------------------
\solution{c10:ex:b1}{\code{app::EchoNode}}
\ans{b} With nothing queued, the stub's first \code{read()} finds its input already exhausted, sets
\code{stop} and returns \code{false}; \code{run()} writes nothing, comes back to the top of its loop,
sees \code{stop}, and returns. The recorded TX is empty.

The mistake only this case catches is \textbf{waiting for input before looking at \code{stop}}: a
\code{run()} that blocks for a first byte, with \code{readBlocking()} or a loop of its own, and only
then starts polling. With three bytes queued the first wait is satisfied at once and the flaw is
invisible; with none, the byte it waits for never comes. Run against both cases, such a
\code{run()} passes the provided one and hangs on the new one until the suite's timeout kills it:

\begin{console}
Test case EchoNode.EchoesQueuedBytesThenStops succeeded!
ERROR: uart_test did not finish within 5s and was killed.
\end{console}

\noindent On the bench that node is deaf to \code{stop} until somebody types. The other natural
mistakes are caught by both cases: echoing without checking what \code{read()} returned sends a byte
that never arrived (the three-byte case records 4 bytes instead of 3, the empty one 1 instead of 0),
and calling \code{readBlocking()} inside the loop hangs both.

One mistake is caught by neither, and it is worth knowing about: a loop that ends when
\code{read()} fails, \code{while (myUart.read(byte))}, instead of when \code{stop} is set. It passes
both cases, because the stub sets \code{stop} at exactly the moment its input runs out, so no test
built on it can tell ``asked to stop'' from ``nothing waiting''. On the bench it returns the first
time the terminal is idle, which is at once.

\ans{c} \code{readBlocking()} spins inside the helper until a byte arrives. While it spins, control
never comes back to \code{run()}'s loop, so \code{stop} is never read again: the node cannot be shut
down, and a test of it hangs, the three-byte case included, the moment its input runs out. The
non-blocking \code{read()} returns on every pass, with a byte or without one, and that return is
what brings the loop back round to look at \code{stop}.

That is also why \code{stop} is a reference read every pass. The flag belongs to someone outside
\code{run()}, and changes while \code{run()} is running: in the test, the stub sets it from inside a
\code{read()} that \code{run()} itself called. Passed by value, \code{run()} would hold a copy made
at the call, \code{false} for ever. Returned, it could not stop a loop that has not returned yet.
Checked once, at the top, it could only stop the loop before it started. A \code{const volatile bool\&} lets \code{run()} see the owner's current value on every pass while promising, through the
\code{const}, never to change it.

And it is \code{volatile} because being read every pass is something the compiler must be told.
\code{run()} never writes the flag, so where nothing the loop calls could write it either --- an
interrupt handler setting it, say --- the compiler is free to load it once and test a register for
ever. In the test the stub happens to set it from inside a \code{read()} the loop calls, which
already forces a fresh load; \code{volatile} makes that true however the flag is set. A byte read
being indivisible on the ATmega is not the same as the compiler being obliged to repeat it.

% ------------------------------------------------------------------------------------------
\solution{c10:ex:c4}{Trace a byte}
One character, typed in the terminal and echoed back:

\begin{booktable}{@{}r>{\raggedright\arraybackslash}p{11.6em}L>{\raggedright\arraybackslash}p{6.4em}@{}}
  \thead{\#} & \thead{Where} & \thead{What happens} & \thead{Built in} \\
  \midrule
  \multicolumn{4}{@{}l}{\itshape In, on the data plane} \\
  1 & Terminal, adapter, \code{rx} pin & A UART frame at 115200~8N1 and 3.3\,V: a start bit, eight
    data bits least significant first, a stop bit & external \\
  2 & \code{sync} & Two flip-flops take the pin into the 50\,MHz domain as \code{rx_s2} &
    \chapref{3}, placed in \chapref{4} \\
  3 & \code{baud_gen} & Ticks at 16 times the baud rate, pacing the receiver & \chapref{2} \\
  4 & \code{uart_rx} & Finds the start edge, samples each bit mid-bit, assembles the byte, pulses
    \code{valid} as \code{rx_push} & \chapref{3} \\
  5 & RX FIFO, a \code{fifo} inside \code{uart_regs} & Stores the byte; \code{STATUS} bit 1 (RX
    valid) becomes 1 & \chapref{4}, \chapref{5} \\
  \midrule
  \multicolumn{4}{@{}l}{\itshape Across the control plane and up the stack} \\
  6 & \code{EchoNode::run()}, \code{Uart::read()} & The node polls; \code{readReg(STATUS)} finds
    RX valid set & \chapref{10}, \chapref{7} \\
  7 & \code{AvrSpi}, the level shifter, \code{spi_slave}, \code{spi_reg_bridge} & Each register
    access is one five-byte transaction; the bridge latches \code{reg_rdata} from the bank's read
    mux & \chapref{9}; provided \\
  8 & \code{readReg(RX_DATA)} & The front byte comes back as \code{00 00 00 XX} on \code{MISO} and
    is assembled most significant first & \chapref{7} \\
  9 & \code{writeReg(RX_POP, 1)} & \code{85 00 00 00 01}; the bridge commits on the fifth byte and
    the FIFO advances & \chapref{7}, \chapref{5} \\
  \midrule
  \multicolumn{4}{@{}l}{\itshape Back down and out} \\
  10 & \code{writeBlocking()}, \code{Uart::write()} & \code{readReg(STATUS)} finds TX ready;
    \code{writeReg(TX_DATA)} sends \code{83 00 00 00 XX} & \chapref{8}, \chapref{7} \\
  11 & \code{uart_regs}, TX FIFO & The committed write pulses \code{tx_push}; the byte is queued &
    \chapref{5}, \chapref{4} \\
  12 & The TX feeder in \code{uart_top} & \code{tx_load} rises with the FIFO non-empty and the
    transmitter idle, starting \code{uart_tx} and popping the FIFO & \chapref{5} \\
  13 & \code{uart_tx}, \code{tx} pin, adapter & The frame \code{'1' & data & '0'}, one bit per 16
    ticks, back to the terminal & \chapref{2} \\
\end{booktable}

The byte changes representation at every boundary it crosses. It is a \textbf{serial frame}, bits
in time, on each UART wire; a \textbf{parallel byte} from \code{uart_rx} onwards, in the RX FIFO and
again in the TX FIFO and the transmitter's frame vector; the \textbf{low byte of a 32-bit register
word} at the register bank, \code{RX_DATA} on the way in and \code{TX_DATA} on the way out;
\textbf{a stream of SPI bytes}, the fourth data byte of a five-byte transaction, on the wire between
the chips; a \code{uint32_t} inside \code{readReg} and \code{writeReg}; and a \code{uint8_t} in the
driver's interface and in \code{EchoNode}.

\textbf{Worth noticing}, now that every layer is in view: one echoed character costs at least five
register transactions, \code{STATUS}, \code{RX_DATA} and \code{RX_POP} to read it, \code{STATUS} and
\code{TX_DATA} to write it back, plus one \code{STATUS} read for every idle pass of \code{run()}. At
about 50\,\textmu s a transaction (Exercise~\ref{c9:ex:3}), that is around 250\,\textmu s per
character, against 86.8\,\textmu s for one frame at 115200~8N1. Typing never notices. A pasted burst
does: the RX FIFO absorbs the difference for a while, but in a continuous burst the backlog passes
its eight entries somewhere around the thirteenth character, and from there bytes are dropped
without a trace, because overrun has no producer in this build: the hardware never sets
\code{ER_OVERRUN}. The bench works; that is its capacity.

